% Practice list D v3 --- topics of Lista 7 (Benigno 2015), QUALIFIER LEVEL (style guide §8: cross-topic,
% parameter analysis, graphs, answer key up to 15 pages).
% \solucoesfalse -> blank list ; \solucoestrue -> answer key + rubric.
% Compile: pdflatex (2x); check: pdffonts (all Type 1). Every number and closed form: practice-L7-check.py
\documentclass[11pt,a4paper]{article}
% Shared preamble of the practice lists (practice-L*.tex), adapted from lista-exam.tex.
% Each list sets \solucoesfalse / \solucoestrue and \listatema right after % Shared preamble of the practice lists (practice-L*.tex), adapted from lista-exam.tex.
% Each list sets \solucoesfalse / \solucoestrue and \listatema right after % Shared preamble of the practice lists (practice-L*.tex), adapted from lista-exam.tex.
% Each list sets \solucoesfalse / \solucoestrue and \listatema right after \input{practice-preamble}.
\usepackage[utf8]{inputenc}
\usepackage[T1]{fontenc}
\usepackage{lmodern}   % Type 1 vector fonts (replaces PK bitmaps)
\usepackage{cmap}      % ToUnicode CMap -> copyable text and accents
\usepackage[english]{babel}
\usepackage{amsmath,amssymb,amsfonts}
\usepackage[a4paper,margin=2.2cm]{geometry}
\usepackage{enumitem}
\usepackage{fancyhdr}
\usepackage{comment}
\usepackage{xcolor}
\usepackage{microtype}
\usepackage{booktabs}
\usepackage{tikz}
\usepackage{pgfplots}
\pgfplotsset{compat=1.16}

\newif\ifsolucoes
\newcommand{\disciplina}{Macroeconomics I --- MPE Insper}
\newcommand{\espacoresol}{4.5cm}   % height of the answer space in the blank list
% The sol environment must be decided after the list sets its toggle.
\AtBeginDocument{\ifsolucoes\includecomment{sol}\else\excludecomment{sol}\fi}

\pagestyle{fancy}\fancyhf{}
\lhead{\footnotesize\disciplina}
\rhead{\footnotesize Practice list\ifsolucoes{} --- Solutions and rubric\fi}
\cfoot{\footnotesize\thepage}
\renewcommand{\headrulewidth}{0.3pt}
\setlength{\headheight}{14pt}
\setlist{topsep=2pt,itemsep=2pt}
\newcounter{questao}

\newcommand{\bloco}[3]{\bigskip\noindent
  {\large\bfseries\color{blue!55!black}Block #1 --- #2}\par\noindent\hfill\textit{(#3)}\par
  \noindent\rule{\linewidth}{0.3pt}\smallskip}
\newcommand{\questao}[1]{\refstepcounter{questao}\medskip\noindent
  \textbf{Question \thequestao.}~#1\par}
\newenvironment{alternativas}{\begin{enumerate}[label=\Alph*.,leftmargin=2.2em,topsep=1pt]}{\end{enumerate}}
\newcommand{\gabarito}[1]{\ifsolucoes\par\noindent\textbf{Answer:}~#1\fi}
\newcommand{\resolespaco}{\ifsolucoes\else\par\textit{Solution:}\vspace{\espacoresol}\par\fi}
\newcommand{\passo}[2]{\par\smallskip\noindent\textbf{Step --- #1.}~#2}
\newcommand{\intuicao}[1]{\par\smallskip\noindent\textit{Economic intuition.}~#1}
\newcommand{\refbib}[1]{\par\smallskip\noindent{\footnotesize\color{black!60}\textbf{Ref:} #1}}
\newcommand{\rubrica}[1]{\par\smallskip\noindent\textbf{Grading rubric.}\begin{itemize}[leftmargin=1.4em,topsep=1pt]#1\end{itemize}}

\newcommand{\cabecalho}{%
\begin{center}
  {\Large\bfseries \disciplina}\\[2pt]
  {\large Practice list: \listatema}\ifsolucoes\\[2pt]\textit{Solutions and grading rubric}\fi
\end{center}
\vspace{2pt}
{\footnotesize\noindent\textbf{Instructions.} Multiple choice: mark the single correct option (no
penalty for wrong answers). \emph{True/False}: every 2 wrong items cancel 1 right one
(score $=\max\{0,\text{right}-\text{wrong}/2\}\times$ item value); blank items do not count.
Discursive questions: show every step. Numbers are illustrative and chosen to give clean answers.\par}}
.
\usepackage[utf8]{inputenc}
\usepackage[T1]{fontenc}
\usepackage{lmodern}   % Type 1 vector fonts (replaces PK bitmaps)
\usepackage{cmap}      % ToUnicode CMap -> copyable text and accents
\usepackage[english]{babel}
\usepackage{amsmath,amssymb,amsfonts}
\usepackage[a4paper,margin=2.2cm]{geometry}
\usepackage{enumitem}
\usepackage{fancyhdr}
\usepackage{comment}
\usepackage{xcolor}
\usepackage{microtype}
\usepackage{booktabs}
\usepackage{tikz}
\usepackage{pgfplots}
\pgfplotsset{compat=1.16}

\newif\ifsolucoes
\newcommand{\disciplina}{Macroeconomics I --- MPE Insper}
\newcommand{\espacoresol}{4.5cm}   % height of the answer space in the blank list
% The sol environment must be decided after the list sets its toggle.
\AtBeginDocument{\ifsolucoes\includecomment{sol}\else\excludecomment{sol}\fi}

\pagestyle{fancy}\fancyhf{}
\lhead{\footnotesize\disciplina}
\rhead{\footnotesize Practice list\ifsolucoes{} --- Solutions and rubric\fi}
\cfoot{\footnotesize\thepage}
\renewcommand{\headrulewidth}{0.3pt}
\setlength{\headheight}{14pt}
\setlist{topsep=2pt,itemsep=2pt}
\newcounter{questao}

\newcommand{\bloco}[3]{\bigskip\noindent
  {\large\bfseries\color{blue!55!black}Block #1 --- #2}\par\noindent\hfill\textit{(#3)}\par
  \noindent\rule{\linewidth}{0.3pt}\smallskip}
\newcommand{\questao}[1]{\refstepcounter{questao}\medskip\noindent
  \textbf{Question \thequestao.}~#1\par}
\newenvironment{alternativas}{\begin{enumerate}[label=\Alph*.,leftmargin=2.2em,topsep=1pt]}{\end{enumerate}}
\newcommand{\gabarito}[1]{\ifsolucoes\par\noindent\textbf{Answer:}~#1\fi}
\newcommand{\resolespaco}{\ifsolucoes\else\par\textit{Solution:}\vspace{\espacoresol}\par\fi}
\newcommand{\passo}[2]{\par\smallskip\noindent\textbf{Step --- #1.}~#2}
\newcommand{\intuicao}[1]{\par\smallskip\noindent\textit{Economic intuition.}~#1}
\newcommand{\refbib}[1]{\par\smallskip\noindent{\footnotesize\color{black!60}\textbf{Ref:} #1}}
\newcommand{\rubrica}[1]{\par\smallskip\noindent\textbf{Grading rubric.}\begin{itemize}[leftmargin=1.4em,topsep=1pt]#1\end{itemize}}

\newcommand{\cabecalho}{%
\begin{center}
  {\Large\bfseries \disciplina}\\[2pt]
  {\large Practice list: \listatema}\ifsolucoes\\[2pt]\textit{Solutions and grading rubric}\fi
\end{center}
\vspace{2pt}
{\footnotesize\noindent\textbf{Instructions.} Multiple choice: mark the single correct option (no
penalty for wrong answers). \emph{True/False}: every 2 wrong items cancel 1 right one
(score $=\max\{0,\text{right}-\text{wrong}/2\}\times$ item value); blank items do not count.
Discursive questions: show every step. Numbers are illustrative and chosen to give clean answers.\par}}
.
\usepackage[utf8]{inputenc}
\usepackage[T1]{fontenc}
\usepackage{lmodern}   % Type 1 vector fonts (replaces PK bitmaps)
\usepackage{cmap}      % ToUnicode CMap -> copyable text and accents
\usepackage[english]{babel}
\usepackage{amsmath,amssymb,amsfonts}
\usepackage[a4paper,margin=2.2cm]{geometry}
\usepackage{enumitem}
\usepackage{fancyhdr}
\usepackage{comment}
\usepackage{xcolor}
\usepackage{microtype}
\usepackage{booktabs}
\usepackage{tikz}
\usepackage{pgfplots}
\pgfplotsset{compat=1.16}

\newif\ifsolucoes
\newcommand{\disciplina}{Macroeconomics I --- MPE Insper}
\newcommand{\espacoresol}{4.5cm}   % height of the answer space in the blank list
% The sol environment must be decided after the list sets its toggle.
\AtBeginDocument{\ifsolucoes\includecomment{sol}\else\excludecomment{sol}\fi}

\pagestyle{fancy}\fancyhf{}
\lhead{\footnotesize\disciplina}
\rhead{\footnotesize Practice list\ifsolucoes{} --- Solutions and rubric\fi}
\cfoot{\footnotesize\thepage}
\renewcommand{\headrulewidth}{0.3pt}
\setlength{\headheight}{14pt}
\setlist{topsep=2pt,itemsep=2pt}
\newcounter{questao}

\newcommand{\bloco}[3]{\bigskip\noindent
  {\large\bfseries\color{blue!55!black}Block #1 --- #2}\par\noindent\hfill\textit{(#3)}\par
  \noindent\rule{\linewidth}{0.3pt}\smallskip}
\newcommand{\questao}[1]{\refstepcounter{questao}\medskip\noindent
  \textbf{Question \thequestao.}~#1\par}
\newenvironment{alternativas}{\begin{enumerate}[label=\Alph*.,leftmargin=2.2em,topsep=1pt]}{\end{enumerate}}
\newcommand{\gabarito}[1]{\ifsolucoes\par\noindent\textbf{Answer:}~#1\fi}
\newcommand{\resolespaco}{\ifsolucoes\else\par\textit{Solution:}\vspace{\espacoresol}\par\fi}
\newcommand{\passo}[2]{\par\smallskip\noindent\textbf{Step --- #1.}~#2}
\newcommand{\intuicao}[1]{\par\smallskip\noindent\textit{Economic intuition.}~#1}
\newcommand{\refbib}[1]{\par\smallskip\noindent{\footnotesize\color{black!60}\textbf{Ref:} #1}}
\newcommand{\rubrica}[1]{\par\smallskip\noindent\textbf{Grading rubric.}\begin{itemize}[leftmargin=1.4em,topsep=1pt]#1\end{itemize}}

\newcommand{\cabecalho}{%
\begin{center}
  {\Large\bfseries \disciplina}\\[2pt]
  {\large Practice list: \listatema}\ifsolucoes\\[2pt]\textit{Solutions and grading rubric}\fi
\end{center}
\vspace{2pt}
{\footnotesize\noindent\textbf{Instructions.} Multiple choice: mark the single correct option (no
penalty for wrong answers). \emph{True/False}: every 2 wrong items cancel 1 right one
(score $=\max\{0,\text{right}-\text{wrong}/2\}\times$ item value); blank items do not count.
Discursive questions: show every step. Numbers are illustrative and chosen to give clean answers.\par}}

\solucoesfalse                 % <<<< \solucoestrue for the ANSWER KEY
\newcommand{\listatema}{The New Keynesian AS--AD Model --- qualifier level (Lista 7)}
\renewcommand{\espacoresol}{6.5cm}

\begin{document}
\cabecalho

{\footnotesize\noindent\textbf{Notation (Benigno 2015).} Lower-case letters are log deviations from the
steady state; bars denote the long run. Households: $u(C)-v(L)+\beta\{u(\bar C)-v(\bar L)\}$,
$u=C^{1-\tilde\sigma^{-1}}/(1-\tilde\sigma^{-1})$, $v=L^{1+\eta}/(1+\eta)$, $\rho\equiv-\ln\beta$. Resources: $y=s_cc+g$,
$\sigma\equiv\tilde\sigma s_c$. Firms: demand $Y(j)=(P(j)/P)^{-\theta}(C+G)$, technology $Y(j)=AL(j)$; a fraction $\alpha$ of
prices is preset at $P^e$. Calibration unless stated: $\alpha=0.75$, $\sigma=0.5$, $\eta=0.2$, $\theta=8$, $\rho=0.01$ (a quarter),
$p^e=\bar p=0$, $g=\bar g=0$, no taxes.\par}

% =============================================================================
\bloco{1}{Foundations: AD, AS, natural and efficient output, taxes}{Benigno \S\S1--5 \,\textbullet\, 25 points}

\questao{(10 points) Objective items.}

\textbf{a)} (3 points) The slope $\kappa$ of the AS curve $p-p^e=\kappa(y-y_n)$ under the calibration is:
\begin{alternativas}
\item $6.600$
\item $0.733$
\item $0.233$
\item $2.200$
\end{alternativas}
\gabarito{B}
\begin{sol}
\passo{derivation}{Flexible firms price at $\tilde p=w-a+\mu$; labor supply gives $w-p=\eta l+\sigma^{-1}(y-g)$ with
$l=y-a$; so $\tilde p-p=(\sigma^{-1}+\eta)(y-y_n)$. With $p=\alpha p^e+(1-\alpha)\tilde p$: $\kappa=\frac{(1-\alpha)(\sigma^{-1}+\eta)}{\alpha}=\frac{0.25\cdot2.2}{0.75}=0.733$.
A inverts $(1-\alpha)/\alpha$; C uses $\sigma+\eta$; D omits the rigidity factor.}
\refbib{Benigno (2015), \S4, eqs.\ (16)--(17), pp.~506--508.}
\end{sol}

\medskip
\textbf{b)} (3 points) \emph{AS--AD $\times$ taxation.} Productivity rises by $a=0.02$ and a payroll tax raises the (log) markup
by $\mu=0.05$, both temporarily. The changes in $(y_n,\ y_e)$ are:
\begin{alternativas}
\item $(-0.0227,\ 0)$
\item $(0.0109,\ 0.0109)$
\item $(-0.0118,\ -0.0118)$
\item $(-0.0118,\ 0.0109)$
\end{alternativas}
\gabarito{D}
\begin{sol}
\passo{coefficients}{$\sigma^{-1}+\eta=2.2$. From (15): $\Delta y_n=(1.2\cdot0.02-0.05)/2.2=-0.0118$; from (19): $\Delta y_e=1.2\cdot0.02/2.2=0.0109$. Taxes
enter only the markup (eq.\ 13), so they move $y_n$, never $y_e$.}
\refbib{Benigno (2015), \S4.4, eqs.\ (13), (15), (19), pp.~507--509.}
\end{sol}

\medskip
\textbf{c)} (2 points, T or F) Because $\theta$ sets the markup, a higher demand elasticity $\theta$ makes the AS curve steeper.
\gabarito{F}
\begin{sol}
\passo{why}{$\theta$ enters $\mu$, hence $y_n$, but not $\kappa=(1-\alpha)(\sigma^{-1}+\eta)/\alpha$: with linear technology and an economy-wide labor
market there is no real rigidity. $\theta$ reappears in the welfare weight of the loss function (Question 4).}
\refbib{Benigno (2015), \S4, eq.\ (17); Romer (2012), \S6.5--6.7.}
\end{sol}

\medskip
\textbf{d)} (2 points, T or F) \emph{AS--AD $\times$ fiscal news.} Holding $i$, $\bar p$ and $p^e$ fixed, news that long-run government
spending will be higher ($\bar g\uparrow$, $g$ unchanged) lowers current output.
\gabarito{T}
\begin{sol}
\passo{channel}{In the AD, $\bar y_n-\bar g$ falls by $\frac{\eta}{\sigma^{-1}+\eta}\Delta\bar g$ (higher future taxes lower long-run consumption);
consumption smoothing lowers current demand; the AS does not move, so $y$ falls by $m_{\bar g}\Delta\bar g$,
$m_{\bar g}=\eta/D=0.067$, $D=(\sigma^{-1}+\eta)(1+\kappa\sigma)$.}
\refbib{Benigno (2015), \S8, Table 1, p.~515.}
\end{sol}

% -----------------------------------------------------------------------------
\questao{(15 points) Discursive. \emph{AS--AD foundations $\times$ taxation $\times$ supply shocks.} Build the model from its
primitives (no taxes unless stated).}

\textbf{a)} (3 points) From the household problem with budget $PC+\frac{\bar P\bar C}{1+i}=WL+\frac{\bar W\bar L}{1+i}+T$, derive the
Euler equation and the labor-supply condition; log-linearize and combine with $y=s_cc+g$ to obtain the AD curve. Give
its slope in the $(y,p)$ plane, explain it, and list its shifters.
\resolespaco
\begin{sol}
\passo{FOCs}{Lagrangian with multiplier $\lambda$: $u_c(C)=\lambda P$, $\beta u_c(\bar C)=\lambda\bar P/(1+i)$, $v_l(L)=\lambda W$; hence
$\frac{u_c(C)}{\beta u_c(\bar C)}=(1+i)\frac P{\bar P}\equiv1+r$ and $\frac{v_l(L)}{u_c(C)}=\frac WP$.}
\passo{log-linear Euler}{$u_c=C^{-1/\tilde\sigma}$: $\bar c-c=\tilde\sigma[i-(\bar p-p)-\rho]$ (with $\ln(1+i)\approx i$, $\rho=-\ln\beta$).}
\passo{AD}{$c=(y-g)/s_c$ and $\bar c=(\bar y-\bar g)/s_c$, and long-run output is $\bar y_n$:
\[y=\bar y_n+(g-\bar g)-\sigma[i-(\bar p-p)-\rho].\]
Slope $dp/dy=-1/\sigma$: a higher $p$ lowers expected inflation $\bar p-p$, raises the real rate and current saving (an
intertemporal channel, no LM curve). Shifters: $i$, $g$, $\bar g$, $\bar p$, $\bar y_n$ (i.e.\ $\bar a$, $\bar\mu$, $\bar g$).}
\refbib{Benigno (2015), \S3, eqs.\ (1)--(6) and (21), pp.~505--509.}
\rubrica{
\item FOCs: 1 pt.
\item Log-linear Euler: 1 pt.
\item AD, slope with its mechanism, shifters: 1 pt. Slope justified by money demand: 0 for this line.
}
\end{sol}

\textbf{b)} (3 points) From the firm's problem derive the optimal flexible price, natural output and the AS curve
with a fraction $\alpha$ of preset prices.
\resolespaco
\begin{sol}
\passo{optimal price}{$\Pi(j)=(P(j)-\frac WA)(\frac{P(j)}P)^{-\theta}(C+G)$; FOC: $(1-\theta)P(j)^{-\theta}+\theta\frac WAP(j)^{-\theta-1}=0\Rightarrow P(j)=\frac{\theta}{\theta-1}\frac WA$;
in logs $\tilde p=w-a+\mu$.}
\passo{gap relation}{Labor supply in logs, $w-p=\eta(y-a)+\sigma^{-1}(y-g)$ (using $l=y-a$, $c=(y-g)/s_c$), gives
$\tilde p-p=(\sigma^{-1}+\eta)y-(1+\eta)a-\sigma^{-1}g+\mu=(\sigma^{-1}+\eta)(y-y_n)$ with $y_n=\frac{(1+\eta)a+\sigma^{-1}g-\mu}{\sigma^{-1}+\eta}$ (where $\tilde p=p$).}
\passo{AS}{$p=\alpha p^e+(1-\alpha)\tilde p\Rightarrow\alpha(p-p^e)=(1-\alpha)(\sigma^{-1}+\eta)(y-y_n)$: $\kappa=\frac{(1-\alpha)(\sigma^{-1}+\eta)}{\alpha}=0.733$.}
\refbib{Benigno (2015), \S4, eqs.\ (9)--(17), pp.~506--508; Romer (2012), problem 6.16.}
\rubrica{
\item Optimal price: 1 pt.
\item Gap relation and $y_n$: 1 pt.
\item AS and $\kappa$: 1 pt. Quoting $\kappa$: 0 for this line.
}
\end{sol}

\textbf{c)} (3 points) Derive efficient output from the planner's problem and compare it with natural output. Show how
labor, payroll and consumption taxes affect each, and find the subsidy that makes the steady state efficient.
\resolespaco
\begin{sol}
\passo{planner}{$\max u(Y-G)-v(Y/A)$: $u_c(Y-G)=v_l(Y/A)/A$; in logs $\sigma^{-1}(y-g)+\eta(y-a)=a$, so $y_e=\frac{(1+\eta)a+\sigma^{-1}g}{\sigma^{-1}+\eta}$ and
$y_n-y_e=-\frac{\mu}{\sigma^{-1}+\eta}$: the markup is a wedge between the real wage and the marginal product of labor.}
\passo{taxes}{With taxes the flexible price becomes $P(j)=\frac{\theta}{\theta-1}\frac{1+\tau_w}{1-\tau_y}\frac WA$ and the labor-supply condition carries
$\frac{1-\tau_l}{1+\tau_c}$; both wedges add up in the aggregate markup $1+\mu=\frac{\theta}{\theta-1}\frac{(1+\tau_w)(1+\tau_c)}{(1-\tau_y)(1-\tau_l)}$ (eq.\ 13). Taxes therefore
lower $y_n$ and leave $y_e$ untouched: the planner faces no tax wedge.}
\passo{efficient subsidy}{$1+\mu=1$ with only a revenue subsidy: $\tau_y=-\frac1{\theta-1}$, i.e.\ 14.3\% at $\theta=8$.}
\refbib{Benigno (2015), \S4.4, eq.\ (19), p.~509; eq.\ (13), p.~507; footnote 21.}
\rubrica{
\item Planner and $y_e$: 1 pt.
\item Taxes in the markup, effect on $y_n$ only: 1.5 pts.
\item Subsidy: 0.5 pt.
}
\end{sol}

\textbf{d)} (3 points) An oil shock raises the markup by $\mu=0.05$ while productivity falls by $a=-0.02$, both for the short run
only. Find the equilibrium output gap in terms of the natural real rate and the policy rate, compute the outcome if
the bank keeps $i=\rho$ and the $i$ that delivers $p=p^e$, and draw the AS--AD diagram (old and new AS, AD, both
equilibria and $y_e$).
\resolespaco
\begin{sol}
\passo{gap}{Substituting the AS into the AD ($g=\bar g$): $(1+\sigma\kappa)(y-y_n)=\bar y_n-y_n-\sigma(i-\rho)+\sigma(\bar p-p^e)$, i.e.
$y-y_n=\frac{\sigma}{1+\sigma\kappa}[r_n-i+(\bar p-p^e)]$ with $r_n\equiv\rho+\sigma^{-1}(\bar y_n-y_n)$.}
\passo{numbers}{$y_n=-0.0336$, $y_e=-0.0109$, $r_n=0.01+2(0.0336)=7.73\%$. With $i=\rho$: gap $=0.366\cdot0.0673=0.0246$, $y=-0.0090$,
$p-p^e=0.0180$; output is above efficient ($y-y_e=+0.0019$). With $i=r_n=7.73\%$: $p=p^e$, $y=y_n$, 2.3\% below $y_e$.}
\begin{center}
\begin{tikzpicture}
\begin{axis}[width=8.6cm,height=6cm,xmin=-4.5,xmax=1.5,ymin=-1.2,ymax=3.4,axis lines=middle,
  xlabel={$y$ ($\times100$)},ylabel={$p$ ($\times100$)},tick label style={font=\footnotesize},label style={font=\small},
  legend style={font=\footnotesize},legend pos=outer north east]
\addplot[thick,black,domain=-1.6:0.6]{-2*x};
\addplot[thick,blue,domain=-1.5:1.5]{0.733333*x};
\addplot[thick,red,domain=-4.5:0]{0.733333*(x+3.3636)};
\addplot[thick,black,dashed,domain=-4.5:-2.8]{-2*(x+3.3636)};
\legend{AD ($i=\rho$),AS,AS$'$,AD$''$ ($i=r_n$)}
\addplot[only marks,mark=*] coordinates {(0,0) (-0.902439,1.804878) (-3.3636,0)};
\addplot[dotted] coordinates {(-1.0909,-1.2) (-1.0909,3.4)};
\node[font=\footnotesize,anchor=south west] at (axis cs:-0.85,1.8) {$E'$};
\node[font=\footnotesize,anchor=south east] at (axis cs:-3.4,0.05) {$E''$};
\node[font=\footnotesize,anchor=north west] at (axis cs:0.05,-0.05) {$E$};
\node[font=\footnotesize,anchor=north east] at (axis cs:-1.12,-0.6) {$y_e$};
\end{axis}
\end{tikzpicture}
\end{center}
\refbib{Benigno (2015), \S5 and \S7, pp.~509--514; natural rate on p.~513.}
\rubrica{
\item Gap formula derived: 1 pt.
\item Numbers for both policies: 1 pt.
\item Graph: AS shifts through $(y_n',p^e)$, AD unchanged at $i=\rho$, $E'$ and $E''$ placed correctly, $y_e$ marked: 1 pt.
}
\end{sol}

\textbf{e)} (3 points) \emph{Parameter analysis.} For the markup shock alone ($\mu=0.05$, $a=0$, $i=\rho$), how do the responses of output and
prices depend on price rigidity? Find them in closed form as functions of $\kappa$, prove the signs, give the limits $\alpha\to0$ and
$\alpha\to1$, fill the table for $\alpha\in\{0.5,0.66,0.75,0.9\}$ and read it.
\resolespaco
\begin{sol}
\passo{closed forms}{$y_n=-\frac{\mu}{\sigma^{-1}+\eta}$, $r_n-\rho=\frac{\sigma^{-1}\mu}{\sigma^{-1}+\eta}$, so the gap is $\frac{\mu}{(1+\sigma\kappa)(\sigma^{-1}+\eta)}$ and
\[y=-\frac{\mu}{\sigma^{-1}+\eta}\cdot\frac{\sigma\kappa}{1+\sigma\kappa},\qquad p-p^e=\frac{\mu}{\sigma^{-1}+\eta}\cdot\frac{\kappa}{1+\sigma\kappa}.\]}
\passo{signs and limits}{$\frac{\partial}{\partial\kappa}\frac{\kappa}{1+\sigma\kappa}=\frac1{(1+\sigma\kappa)^2}>0$: more flexibility (lower $\alpha$, higher $\kappa$) means a larger fall in output and
a larger rise in prices. $\alpha\to1$ ($\kappa\to0$): $y\to0$, $p\to p^e$ (output is demand-determined). $\alpha\to0$ ($\kappa\to\infty$): $y\to y_n$ and
$p-p^e\to\frac{\mu}{1+\sigma\eta}=4.5\%$.}
\passo{table}{
\begin{center}
\begin{tabular}{lcccc}
\toprule
$\alpha$ & 0.5 & 0.66 & 0.75 & 0.9 \\
\midrule
$\kappa$ & 2.200 & 1.133 & 0.733 & 0.244 \\
$y$ (\%) & $-1.19$ & $-0.82$ & $-0.61$ & $-0.25$ \\
$p-p^e$ (\%) & 2.38 & 1.64 & 1.22 & 0.50 \\
\bottomrule
\end{tabular}
\end{center}}
\passo{reading}{Rigid prices shield output from a cost shock when the policy rate is held fixed, because demand pins
output down; flexible prices pass the shock to prices and output falls toward $y_n$. With the markup shock alone
$y_e$ does not move, so every response leaves output below the efficient level, and the more flexible the prices, the
larger both the efficient gap and the price surprise (in (d), output could exceed $y_e$ only because $y_e$ fell too).}
\refbib{Benigno (2015), \S7, p.~513; Romer (2012), problem 6.16.}
\rubrica{
\item Closed forms: 1 pt.
\item Signs and both limits proved: 1 pt.
\item Table and reading: 1 pt.
}
\end{sol}

% =============================================================================
\bloco{2}{Policy: optimal rules, the zero bound and household debt}{Benigno \S\S6--12 \,\textbullet\, 25 points}

\questao{(10 points) Objective items.}

\textbf{a)} (3 points) \emph{Monetary $\times$ fiscal policy.} A temporary cut in consumption taxes of 2 points (like the 2009 IPI
cut on cars, modeled as $\tau_c=-0.02$, entering the markup one for one) changes the natural real rate
$r_n=\rho+\sigma^{-1}(\bar y_n-y_n)+\sigma^{-1}(g-\bar g)+(\bar\tau_c-\tau_c)$ by:
\begin{alternativas}
\item $+2.00$ points
\item $-1.82$ points
\item $+0.18$ points
\item $0$
\end{alternativas}
\gabarito{C}
\begin{sol}
\passo{two channels}{Supply: $\Delta\mu=-0.02$, $\Delta y_n=0.02/2.2=0.0091$, lowering $r_n$ by $0.0182$. Demand: $\bar\tau_c-\tau_c$ rises by $0.02$. Net
$+0.0018$: the bank must raise $i$ slightly to keep output at its natural level.}
\refbib{Benigno (2015), \S8 and p.~513.}
\end{sol}

\medskip
\textbf{b)} (3 points) \emph{Fiscal multipliers.} With $i$ fixed, a temporary rise in $g$ of 1\% of GDP changes output and the gap
$y-y_n$ by, respectively:
\begin{alternativas}
\item $1.000\%$ and $0.000\%$
\item $0.909\%$ and $0.000\%$
\item $0.732\%$ and $0.067\%$
\item $0.976\%$ and $0.067\%$
\end{alternativas}
\gabarito{D}
\begin{sol}
\passo{both curves move}{$\Delta y_n=\frac{\sigma^{-1}}{\sigma^{-1}+\eta}=0.909$ per unit; $\Delta r_n=\sigma^{-1}(1-0.909)=0.182$; gap $=0.366\cdot0.182=0.067$; output
$0.909+0.067=0.976=m_g$ (Table 1).}
\refbib{Benigno (2015), \S8, Table 1, p.~515.}
\end{sol}

\medskip
\textbf{c)} (2 points, T or F) At the zero lower bound with $r_n<0$, a credible commitment to a long-run price level higher
by $-r_n$ restores $y=y_n$ and $p=p^e$.
\gabarito{T}
\begin{sol}
\passo{gap formula}{With $i=0$, $y-y_n=\frac{\sigma}{1+\sigma\kappa}[r_n+(\bar p-p^e)]=0$ when $\bar p-p^e=-r_n$; then $p=p^e$.}
\refbib{Benigno (2015), \S9, p.~517.}
\end{sol}

\medskip
\textbf{d)} (2 points, T or F) \emph{AS--AD $\times$ household debt.} In the borrower--saver economy with constant inflation between
the short and the long run, more price flexibility reduces the output cost of a deleveraging shock.
\gabarito{F}
\begin{sol}
\passo{paradox of flexibility}{The AD slopes up (Fisher debt deflation) and the output response $\frac{(1-\chi)\hat d}{\chi-(1-\chi)d_0\kappa}$ grows in size with $\kappa$
(Question 4(e)).}
\refbib{Benigno (2015), \S10, pp.~519--521.}
\end{sol}

% -----------------------------------------------------------------------------
\questao{(15 points) Discursive. \emph{Optimal policy $\times$ zero lower bound $\times$ household debt.}}

\textbf{a)} (3 points) The bank minimizes Benigno's welfare loss $L=\tfrac12(y-y_e)^2+\tfrac12\tfrac{\theta}{\kappa}(p-p^e)^2$ (eq.\ 33) subject to the AS. Find
the optimal relation between the efficient gap and price surprises, describe it in the $(y,p)$ plane and prove that shocks
that leave $\mu$ unchanged create no trade-off.
\resolespaco
\begin{sol}
\passo{substitution}{With $p-p^e=\kappa(y-y_n)$: $L=\tfrac12(y-y_e)^2+\tfrac12\theta\kappa(y-y_n)^2$, strictly convex in $y$.}
\passo{targeting rule}{FOC: $(y-y_e)+\theta\kappa(y-y_n)=0$, i.e.\ $(y-y_e)+\theta(p-p^e)=0$ (eq.\ 34); $y^*=\frac{y_e+\theta\kappa y_n}{1+\theta\kappa}$. In the plane: a line through
$(y_e,p^e)$ with slope $-1/\theta$; the optimum is its intersection with the AS.}
\passo{no trade-off}{If $\mu$ is unchanged, $y_n=y_e$, and $y=y_e$, $p=p^e$ satisfies both the AS and the rule with $L=0$ (divine
coincidence); it is implemented by $i=r_n+(\bar p-p^e)$.}
\refbib{Benigno (2015), \S11, eqs.\ (33)--(34), pp.~521--522.}
\rubrica{
\item Substitution and FOC: 1.5 pts.
\item Geometry: 0.5 pt.
\item Divine coincidence with the argument: 1 pt.
}
\end{sol}

\textbf{b)} (3 points) For the shock of Question 2(d), compute the optimal $y$, $p-p^e$ and $i$ and compare with strict price
stability. Then show that no Taylor rule $i=\bar\imath+\psi_y(y-y_e)+\psi_\pi(p-p^e)$ with a constant intercept implements the optimum under
productivity shocks ($a=\pm0.004$).
\resolespaco
\begin{sol}
\passo{optimum}{$\theta\kappa=5.867$: $y^*=\frac{-0.0109+5.867(-0.0336)}{6.867}=-0.0303$, $p-p^e=\kappa(y^*-y_n)=0.0024$; check
$(y^*-y_e)+\theta(p-p^e)=-0.0194+0.0194=0$.}
\passo{interest rate}{$i=r_n-\frac{1+\sigma\kappa}{\sigma}(y^*-y_n)=0.0773-2.733\cdot0.0033=6.82\%$, between the passive 1\% and strict price
stability's 7.73\%: the bank accepts a 0.24\% price rise to cushion output.}
\passo{Taylor rule}{Under a pure productivity shock $y_n=y_e$ and the optimum has zero gaps, so any Taylor rule prescribes
$i=\bar\imath$. The optimum needs $i=r_n$: $a=0.004$ gives $r_n=0.56\%$, $a=-0.004$ gives $1.44\%$. One intercept cannot equal both:
the rule must track the natural rate.}
\refbib{Benigno (2015), \S11.}
\rubrica{
\item Optimal $y$, $p$ with the check: 1 pt.
\item Optimal $i$ and comparison: 1 pt.
\item Taylor-rule argument: 1 pt.
}
\end{sol}

\textbf{c)} (3 points) Pessimism: long-run productivity is expected to fall, $\bar a=-0.05$, with $a=0$. Show that the zero bound binds,
compute output and prices at $i=0$ with $\bar p=p^e$, find the $\bar p$ that restores the optimum, give the limits as $\kappa\to\infty$ and
$\kappa\to0$, and discuss credibility.
\resolespaco
\begin{sol}
\passo{natural rate}{$\bar y_n=\frac{1.2}{2.2}(-0.05)=-0.0273$, $y_n=y_e=0$, $r_n=0.01+2(-0.0273)=-4.45\%<0$: $i=r_n$ is unavailable.}
\passo{at $i=0$}{Gap $=0.366\cdot(-0.0445)=-0.0163$; $p-p^e=-0.0120$. Remedy: $\bar p-p^e=-r_n=4.45\%$.}
\passo{limits and credibility}{$\kappa\to\infty$: gap $\to0$ and $p-p^e\to r_n=-4.45\%$; $\kappa\to0$: $p=p^e$ and gap $=\sigma r_n=-0.0223$. Once the recession
is over the bank prefers $\bar p=p^e$: the promise is time-inconsistent and works only if believed.}
\refbib{Benigno (2015), \S6.3 and \S9, pp.~513 and 517.}
\rubrica{
\item $r_n$ and the binding bound: 1 pt.
\item Outcome at $i=0$ and the required $\bar p$: 1 pt.
\item Limits and credibility: 1 pt.
}
\end{sol}

\textbf{d)} (3 points) With borrowers and savers (share $\chi$ of savers, initial debt $d_0$), Benigno's AD is
\[y=\bar y_n-\varphi[i-(\bar p_n-p)-\rho]+\tfrac1\chi(g-\bar g)+\tfrac{1-\chi}{\chi}[\hat d+d_0(p-p^e)],\qquad\varphi\equiv\frac{\sigma\chi+(1-\chi)d_0\beta}{\chi},\]
where $\hat d$ is the change in the debt limit. Find $dy/dp$ when $\bar p_n$ is anchored and when inflation $\bar p_n-p$ is constant,
name the two channels, and derive the spending multipliers in both regimes; compute them with $\chi=2/3$, $d_0=1.2$, $\beta=0.98$.
\resolespaco
\begin{sol}
\passo{slopes}{Anchored: $dy/dp=-\varphi+\frac{(1-\chi)d_0}{\chi}=-\varpi$, $\varpi\equiv\sigma-\frac{d_0(1-\beta)(1-\chi)}{\chi}=0.488$: the real-rate channel outweighs the
Fisher channel (a higher $p$ erodes nominal debt). Constant inflation: the real-rate term no longer depends on $p$, so
$dy/dp=+\frac{(1-\chi)d_0}{\chi}>0$: the AD slopes up.}
\passo{multipliers}{Temporary $g$ moves $y_n$ by $\frac{g}{1+\sigma\eta}$ and the AD by $g/\chi$. Anchored: $m=\frac{1/\chi+\varpi\kappa/(1+\sigma\eta)}{1+\varpi\kappa}$. Constant inflation:
$m=\frac{1-(1-\chi)d_0\kappa/(1+\sigma\eta)}{\chi-(1-\chi)d_0\kappa}$, valid when $\chi>(1-\chi)d_0\kappa$.}
\passo{numbers}{$\alpha=0.75$: $1.344$ and $1.964$; $\alpha=0.66$: $1.290$ and $2.756$ (Benigno's figures). Both exceed one: Ricardian equivalence
fails for constrained borrowers.}
\refbib{Benigno (2015), \S10, eq.\ (32), pp.~519--521.}
\rubrica{
\item Both slopes with the two channels: 1 pt.
\item Both multipliers derived: 1.5 pts. Quoted from the article: max.\ 0.5 pt.
\item Numbers: 0.5 pt.
}
\end{sol}

\textbf{e)} (3 points) \emph{Parameter analysis.} How does the output response to a fall in the debt limit $\hat d=-0.1$ depend on price
flexibility in each regime? Find it in closed form, prove the signs, give the stability bound, fill the table for
$\alpha\in\{0.66,0.75,0.9\}$, and draw both responses against $\kappa$.
\resolespaco
\begin{sol}
\passo{closed forms and signs}{Anchored: $y=\frac{(1-\chi)\hat d/\chi}{1+\varpi\kappa}$, smaller in size as $\kappa$ rises. Constant inflation: $y=\frac{(1-\chi)\hat d}{\chi-(1-\chi)d_0\kappa}$, larger in
size as $\kappa$ rises, and unbounded as $\kappa\to\frac{\chi}{(1-\chi)d_0}=1.667$, the stability bound.}
\passo{table}{
\begin{center}
\begin{tabular}{lccc}
\toprule
$\alpha$ ($\kappa$) & 0.66 (1.133) & 0.75 (0.733) & 0.90 (0.244) \\
\midrule
$y$ anchored & $-3.2\%$ & $-3.7\%$ & $-4.5\%$ \\
$y$ constant inflation & $-15.6\%$ & $-8.9\%$ & $-5.9\%$ \\
multiplier, anchored & 1.29 & 1.34 & 1.44 \\
multiplier, constant inflation & 2.76 & 1.96 & 1.60 \\
\bottomrule
\end{tabular}
\end{center}}
\begin{center}
\begin{tikzpicture}
\begin{axis}[width=10cm,height=5.4cm,xmin=0,xmax=1.6,ymin=-30,ymax=0,xlabel={$\kappa$},ylabel={$y$ (\%)},
  legend style={font=\footnotesize},legend pos=south west,tick label style={font=\footnotesize},label style={font=\small}]
\addplot[thick,blue,domain=0:1.6,samples=80]{100*(-0.05)/(1+0.488*x)};
\addplot[thick,red,domain=0:1.55,samples=120]{100*(-0.1/3)/(2/3-0.4*x)};
\addplot[dotted] coordinates {(1.6667,-30) (1.6667,0)};
\legend{anchored $\bar p_n$,constant inflation}
\end{axis}
\end{tikzpicture}
\end{center}
\passo{reading}{When long-run prices are anchored, flexible prices help: a lower $p$ raises expected inflation and cuts the
real rate. When inflation is expected to persist, a lower $p$ lowers future prices too and raises the real value of
debt: flexibility deepens the slump (paradox of flexibility), and the multiplier is largest exactly where the slump is.}
\refbib{Benigno (2015), \S10, pp.~519--521, and footnote 20.}
\rubrica{
\item Closed forms, signs and the bound: 1 pt.
\item Table: 1 pt.
\item Graph: one decreasing-in-size and one exploding curve, bound marked: 1 pt.
}
\end{sol}

\end{document}
